\section{Conclusions}
\label{sec:conclusions}

On synthetic single-hop mother-of retrieval, RI poorly identifies causally
important heads: it misses major contributors and selects many with weak
standalone effects (Section~\ref{sec:results-why}). Path interventions reveal
interacting fact-position writers and answer-position readers. RI-selected
heads transmit fact-dependent information through value channels and influence
later readers' queries; copying alone does not explain these roles
(Sections~\ref{sec:results-char}--\ref{sec:results-routes};
Appendix~\ref{app:stage3}). This establishes task-specific participation,
not that individual heads encode the relation.

Weaker RI heads contribute collectively, but their signed group effect
depends on the replacement baseline (Section~\ref{sec:results-routes}).
L9H16 and L1H27 reproduce binding gains on held-out families, though their
fact-swap effects nearly vanish with paired donors
(Appendix~\ref{app:s44-validation}). The retained 50-head set does not
preserve the full task, and these weaker heads' routes remain unresolved
(Appendix~\ref{app:s44-results}). Early routing changes during training do
not identify mature high-RI heads (Appendix~\ref{app:dev}). Thus, neither RI
nor training-time association identifies a sufficient causal circuit.
